% ==========================================
% CHAPTER 2
% ==========================================
\chapter{Basic Python} % 

\section{Introduction} % 
Programming is the type of thing whose uses and applicabilities seem extremely straightforward once you know how to do it, and extremely nebulous and intimidating before that point. %  

\textbf{Definition 2.0.1} Python is a programming language (and yes, it's named for the sketch troupe Monty Python). It is an interpreted high-level programming language for general-purpose programming, and one of the most common languages used in astronomy. % 

\section{Hello, World!} % 
In Python, this is dead easy. From the terminal, simply type: % 

\begin{lstlisting}
ipython
\end{lstlisting} % 

Now, all you have to do is type: % 

\begin{lstlisting}
print('hello world!')
\end{lstlisting} % 

\textbf{Definition 2.1.1} A print statement is a line of code which tells the interpreter to output something to the screen. % 

\section{Data types} % 
Python, like most programming languages, divides up all the possible ``things'' you can play around with into what are called data types. % 

Some basic data types are listed and defined below: % 
\begin{itemize}
    \item \textbf{Integers:} The counting numbers. Ex: -1, 0, 1, 2, 3... % 
    \item \textbf{Floats:} Decimal numbers. Ex: 1.0, 2.345, 6.3... % 
    \item \textbf{Strings:} An iterable data type most commonly used to hold words/phrases. Ex: \texttt{"cat"}, \texttt{"/home/ipasha"} % 
    \item \textbf{Lists:} Stored lists of any combination of data types. Ex: \texttt{[1, 2, 'star', 'fish']} % 
    \item \textbf{Dictionaries:} A collection of pairs, where one is a ``key'' and the other is a ``value.'' % 
\end{itemize}

\section{Basic Math} % 
Within the python interpreter you can perform mathematical operations: % 

\begin{lstlisting}
[IN]: 3 + 5
[OUT]: 8
\end{lstlisting} % 

\section{Variables} % 
\textbf{Definition 2.4.1} A variable is a user-defined, symbolic name which points to a spot in a computer's memory where a value has been stored. % 

Declaring variables in Python is easy: % 

\begin{lstlisting}
[IN]: x = 5.0
[IN]: y = 'cat'
\end{lstlisting} % 

\section{Storing and Manipulating Data in Python} % 
\subsection{Arrays vs Lists} % 
Lists and numpy arrays are the natural data type with which to store data. %  Let's define a list to play with: % 

\begin{lstlisting}
[IN]: my_list = [1,5,2,7,3,7,8]
\end{lstlisting} % 

If we want to multiply all the distances by 2, numpy arrays make this much easier and faster than standard Python lists: % 

\begin{lstlisting}
[IN]: my_array = np.array([1,5,2,7,3,7,8])
[IN]: my_array * 2
[OUT]: array([2,10,4,14,6,14,16])
\end{lstlisting} % 

\section{Array, String, and List Indexing} % 
The procedure by which we extract subsets of a larger array/list is known as slicing, or indexing. %  
By convention, this index starts with 0. % 

\begin{lstlisting}
[IN]: list_1 = [1, 2, 4, 'cat', 5]
[IN]: x = list_1[0]
[IN]: print(x)
[OUT]: 1
\end{lstlisting} % 

\section{Dictionaries} % 
\textbf{Definition 2.8.1} A dictionary is a Python container, like a list or array, but which uses ``keys'' instead of indices to specify elements. % 

\begin{lstlisting}
simple_dict = {'key1': value1, 'key2': value2}
pulled_value = simple_dict['key1']
\end{lstlisting} % 